\documentclass[11pt]{article}
\usepackage[margin=1in]{geometry}
\usepackage{amsmath,amssymb,amsthm,mathtools}
\usepackage{booktabs}
\usepackage[colorlinks=true,linkcolor=blue!55!black,citecolor=blue!55!black,urlcolor=blue!55!black]{hyperref}
\usepackage{xcolor}
\usepackage{microtype}

\newtheorem{theorem}{Theorem}
\newtheorem{proposition}[theorem]{Proposition}
\newtheorem{lemma}[theorem]{Lemma}
\newcommand{\eps}{\{-1,1\}}
\newcommand{\E}{\mathbb E}

\hypersetup{
  pdftitle={An AI Attempt at a Min-Max Problem for Quadratic Forms of Signs},
  pdfauthor={OpenAI Codex},
  pdfsubject={Partial results for MathOverflow question 413935},
  pdfkeywords={quadratic forms, signs, discrepancy, conference matrices, cut codes}
}

\title{\textbf{An AI Attempt at a Min--Max Problem for\\
Quadratic Forms of Signs}\\[4pt]
\large Exact reformulations and the bounds \(1/\pi\) and \(1/2\)}
\author{OpenAI Codex\\
\small AI-generated research note}
\date{25 July 2026}

\begin{document}
\maketitle

\begin{center}
\fcolorbox{red!65!black}{red!4}{
\begin{minipage}{0.89\textwidth}
\textbf{Status and nonclaim.}
This note does not prove that the limit in Question~\ref{q:main} exists,
and it does not prove that the limit fails to exist. It is an unsuccessful
solution attempt containing several rigorous partial results. It should not
be submitted or cited as a solution to the MathOverflow problem.
\end{minipage}}
\end{center}

\begin{abstract}
For
\[
 F(n)=\min_{a_{ij}\in\eps}
 \max_{x\in\eps^n}
 \left|\sum_{1\le i<j\le n}a_{ij}x_ix_j\right|,
\]
we give an exact covering-radius formulation and prove
\[
 \frac1\pi
 \le \liminf_{n\to\infty}\frac{F(n)}{n^{3/2}}
 \le \limsup_{n\to\infty}\frac{F(n)}{n^{3/2}}
 \le \frac12.
\]
The lower bound follows from two correlated Gaussian sign-roundings and
the arcsine law. The upper bound follows from Paley conference matrices
and principal submatrices. We also record why monotonicity and ordinary
block composition do not settle convergence. The question whether the
normalized sequence converges remains open in this attempt.
\end{abstract}

\section{The question}

\label{q:main}
The problem, posed on MathOverflow in 2022 \cite{mo}, asks whether
\[
 \lim_{n\to\infty}\frac{F(n)}{n^{3/2}}
\]
exists. The emphasis is on existence, not merely on proving the
\(\Theta(n^{3/2})\) scale.

Given the coefficients \(a_{ij}\), let \(A\) be the symmetric matrix with
zero diagonal and off-diagonal entries \(A_{ij}=a_{ij}\). Write
\[
 Q_A(x)=\sum_{i<j}a_{ij}x_ix_j=\frac12x^{\mathsf T}Ax,
 \qquad
 M(A)=\max_{x\in\eps^n}|Q_A(x)|.
\]
Then \(F(n)=\min_A M(A)\).

\section{Switching and an exact code formulation}

Changing \(x\) conjugates \(A\) by a diagonal sign matrix. Thus \(M(A)\)
is the largest absolute total signed edge sum in the switching class of
\(A\). The relation between switching classes and cosets of the cutset
code is standard \cite{sole-zaslavsky}; the absolute-value version needed
here has the following exact form.

\begin{proposition}[Augmented cut-code identity]
\label{prop:code}
Put \(m=\binom n2\) and
\[
 c_x=(x_ix_j)_{i<j}\in\eps^m,\qquad
 D_n=\{c_x:x\in\eps^n\}\cup\{-c_x:x\in\eps^n\}.
\]
If \(\rho(D_n)\) denotes the covering radius of \(D_n\) in the Hamming
cube \(\eps^m\), then
\[
 \boxed{F(n)=m-2\rho(D_n).}
\]
\end{proposition}

\begin{proof}
For \(a\in\eps^m\),
\[
 \langle a,c_x\rangle=m-2d_H(a,c_x).
\]
Because \(D_n\) contains \(c_x\) and \(-c_x\),
\[
 \max_x|\langle a,c_x\rangle|
 =m-2d_H(a,D_n).
\]
Minimizing the left side over \(a\) maximizes the final distance, giving
the stated formula.
\end{proof}

In particular, the original question asks whether the second-order
covering-radius deficit
\[
 \frac m2-\rho(D_n)=\frac{F(n)}2
\]
has an \(n^{3/2}\)-normalized limit.

\section{A Gaussian-sign lower bound}

The next argument is uniform over every admissible coefficient matrix.

\begin{theorem}
\label{thm:lower}
For every \(n\ge2\),
\[
 \boxed{F(n)\ge \frac{n\sqrt{n-1}}{\pi}.}
\]
Consequently,
\[
 \liminf_{n\to\infty}\frac{F(n)}{n^{3/2}}\ge\frac1\pi.
\]
\end{theorem}

\begin{proof}
Fix \(A\), let \(g\sim N(0,I_n)\), and choose \(t>0\). Define
\[
 z^\pm=\left(I\pm\frac{t}{\sqrt n}A\right)g,
 \qquad x_i^\pm=\operatorname{sgn}(z_i^\pm).
\]
Every coordinate of \(z^\pm\) has variance
\[
 d=1+t^2\frac{n-1}{n}.
\]
For \(i\ne j\),
\[
 \operatorname{Cov}(z_i^\pm,z_j^\pm)
 =\pm\frac{2t}{\sqrt n}a_{ij}
   +\frac{t^2}{n}(A^2)_{ij}.
\]
Let \(r_{ij}^\pm\) be the corresponding correlation and set
\[
 u_{ij}=\frac{a_{ij}t^2(A^2)_{ij}}{nd},
 \qquad
 v=\frac{2t}{\sqrt n\,d}.
\]
Then \(a_{ij}r_{ij}^\pm=u_{ij}\pm v\).

For centered jointly Gaussian variables with correlation \(r\), the
arcsine law states
\[
 \E[\operatorname{sgn}(Z_1)\operatorname{sgn}(Z_2)]
 =\frac2\pi\arcsin r.
\]
For completeness, write
\((Z_1,Z_2)=(G_1,rG_1+\sqrt{1-r^2}G_2)\) with independent standard
Gaussians. The direction of \((G_1,G_2)\) is uniform on the circle, and
the angular measure of the two same-sign sectors gives the displayed
identity.
Since \(\arcsin\) is odd and its derivative is at least \(1\) on
\((-1,1)\),
\[
\begin{aligned}
a_{ij}\left(\arcsin r_{ij}^+-\arcsin r_{ij}^-\right)
 &=\arcsin(u_{ij}+v)-\arcsin(u_{ij}-v)\\
 &\ge 2v.
\end{aligned}
\]
The endpoints are valid correlations, so they lie in \([-1,1]\); the
same inequality follows at an endpoint by continuity.

Summing over \(m=\binom n2\) edges gives
\[
\begin{aligned}
\E Q_A(x^+)-\E Q_A(x^-)
 &=\frac2\pi\sum_{i<j}a_{ij}
   \left(\arcsin r_{ij}^+-\arcsin r_{ij}^-\right)\\
 &\ge\frac{4mv}{\pi}
 =\frac{8tm}{\pi\sqrt n\,d}.
\end{aligned}
\]
Both expectations belong to \([-M(A),M(A)]\). Therefore
\[
 M(A)\ge\frac{4tm}{\pi\sqrt n\,d}.
\]
The ratio \(t/(1+t^2(n-1)/n)\) is maximized at
\(t=\sqrt{n/(n-1)}\), where \(d=2\). Substitution yields
\[
 M(A)\ge \frac{2m}{\pi\sqrt{n-1}}
 =\frac{n\sqrt{n-1}}{\pi}.
\]
The bound holds for every \(A\), completing the proof.
\end{proof}

\paragraph{Independent check.}
A separate finite verifier reconstructs the covariance entries for
random sign matrices, checks the arcsine inequality edge by edge, checks
the optimized algebra through order \(9999\), and includes corruption
controls. This is a check of the displayed calculation, not a substitute
for the proof.

\section{A conference-matrix upper bound}

\begin{theorem}
\label{thm:upper}
\[
 \boxed{\limsup_{n\to\infty}\frac{F(n)}{n^{3/2}}\le\frac12.}
\]
\end{theorem}

\begin{proof}
A symmetric conference matrix \(C\) of order \(N\) has zero diagonal,
off-diagonal entries in \(\eps\), and
\[
 C^2=(N-1)I.
\]
For every \(x\in\eps^N\),
\[
 |Q_C(x)|
 =\frac12|x^{\mathsf T}Cx|
 \le\frac12\|x\|_2\|Cx\|_2
 =\frac12N\sqrt{N-1}.
\]

The Paley construction supplies such a matrix of order \(q+1\) whenever
\(q\) is a prime power with \(q\equiv1\pmod4\); see, for example,
\cite{cameron}. The prime number theorem in arithmetic progressions
allows, for each sufficiently large \(n\), a prime
\(q\equiv1\pmod4\) with \(q+1\ge n\) and \((q+1)/n\to1\).

Take any \(n\)-by-\(n\) principal submatrix \(A\) of the conference
matrix of order \(q+1\). Compression cannot increase operator norm, so
\(\|A\|_{\mathrm{op}}\le\sqrt q\). Hence
\[
 F(n)\le M(A)
 \le\frac12n\sqrt q
 =\left(\frac12+o(1)\right)n^{3/2}.
\]
\end{proof}

\section{Elementary regularity and the composition wall}

\begin{proposition}
\label{prop:monotone}
For every \(n\ge2\),
\[
 F(n)\le F(n+1)\le F(n)+n.
\]
\end{proposition}

\begin{proof}
Restrict a signing on \(n+1\) vertices to its first \(n\) vertices.
For fixed \(x\in\eps^n\), its restricted energy is the average of the
two full energies obtained by setting the final sign to \(1\) and
\(-1\). The restricted norm is therefore no larger than the full norm,
which proves the first inequality after minimization.

For the second inequality, extend an optimal \(n\)-vertex signing by
arbitrary incident signs \(b_i\). If \(L(x)=\sum_i b_ix_i\), then
\[
 \max_{s\in\eps}|Q_A(x)+sL(x)|=|Q_A(x)|+|L(x)|
 \le F(n)+n.
\]
\end{proof}

This regularity is too weak to force convergence. A direct block
construction also loses at leading order. If optimal arrays of orders
\(n\) and \(k\) are joined by independently chosen cross-edge signs,
a union bound gives a cross discrepancy of order
\[
 \sqrt{2(\log2)\,nk(n+k)}.
\]
When \(n\) and \(k\) are comparable, this is of order
\((n+k)^{3/2}\), exactly the target scale. Thus neither Fekete's lemma
nor ordinary approximate subadditivity applies.

The analogous rectangular switching game is also governed by
unimodular bilinear forms and has a long discrepancy literature
\cite{brown-spencer,pellegrino-raposo}. Those bounds do not provide the
lossless symmetric composition needed here.

\section{What remains}

Combining Theorems~\ref{thm:lower} and~\ref{thm:upper} gives the proved
partial result
\[
\boxed{
\frac1\pi
\le \liminf_{n\to\infty}\frac{F(n)}{n^{3/2}}
\le \limsup_{n\to\infty}\frac{F(n)}{n^{3/2}}
\le\frac12.
}
\]
Proposition~\ref{prop:code} gives an exact second description:
\[
\boxed{F(n)=\binom n2-2\rho(D_n).}
\]

The first missing implication is still a limit theorem. One needs either
an asymptotically lossless amplification taking a near-optimal signing
at one order to a dense set of larger orders, or a different
interpolation that keeps the cross-block contribution at the same
optimized constant. The direct block, tensor, and annealed
partition-function approaches examined in this attempt all incur a
leading-order loss.

\medskip
\noindent\textbf{Final status.}
This note narrows the possible asymptotic interval and supplies an exact
coding formulation. It does not answer the question posed on
MathOverflow.

\begin{thebibliography}{9}

\bibitem{mo}
P.~Ivanisvili,
\emph{Min-max of a quadratic form of plus-minus ones},
MathOverflow question 413935, 2022.
\url{https://mathoverflow.net/questions/413935/}

\bibitem{sole-zaslavsky}
P.~Sol\'e and T.~Zaslavsky,
\emph{A coding approach to signed graphs},
SIAM J. Discrete Math. \textbf{7} (1994), 544--553.
\url{https://doi.org/10.1137/S0895480189174374}

\bibitem{cameron}
P.~J. Cameron,
\emph{Conference matrices}, lecture notes.
\url{https://webspace.maths.qmul.ac.uk/p.j.cameron/csgnotes/conftalk.pdf}

\bibitem{brown-spencer}
T.~A. Brown and J.~H. Spencer,
\emph{Minimization of \(\pm1\) matrices under line shifts},
Colloq. Math. \textbf{23} (1971), 165--171.

\bibitem{pellegrino-raposo}
D.~Pellegrino and A.~Raposo Jr.,
\emph{Upper bounds for the constants of Bennett's inequality and the
Gale--Berlekamp switching game},
arXiv:2111.00445, 2022.
\url{https://arxiv.org/abs/2111.00445}

\end{thebibliography}

\end{document}
